\begin{figure}[tbp]
\centering
\begin{tikzpicture}[x=0.94cm,y=0.94cm,font=\small]
% 顶部：中断源
\node[box,text width=2.8cm] (src) at (9.6,3.4) {Interrupt Source};
% 第二层：Device table 与 Interrupt translation table
\node[greenbox,text width=2.6cm] (devt) at (3.2,1.7) {Device table};
\node[greenbox,text width=3.6cm] (itt) at (9.2,1.7) {Interrupt translation\\table};
% 第三层：Collection table（物理）与 vPE table（GICv4）
\node[box,text width=2.6cm] (ct) at (3.8,-1.3) {Collection table};
\node[amberbox,text width=3.2cm] (vpet) at (13.6,-1.3) {vPE table\\(GICv4 only)};
% 第四层：两个 Redistributor 与各自的 pending 表
\node[greenbox,text width=2.8cm] (vpend) at (13.6,-4.0) {Virtual\\pending table};
\node[greenbox,text width=2.6cm] (pend) at (3.8,-4.0) {Pending table};
\node[note,font=\footnotesize\bfseries] (rda) at (13.6,-5.1) {Redistributor A};
\node[note,font=\footnotesize\bfseries] (rdb) at (3.8,-5.1) {Redistributor B};
% 连线与 ID 标注
\draw[wire] (src.south west) -- node[above left,font=\scriptsize,pos=0.45]{DeviceID} (devt.north east);
\draw[wire] (devt.east) -- node[below,font=\scriptsize,pos=0.5]{ITT base} (itt.west);
\draw[wire] (src.south) -- node[left,font=\scriptsize,pos=0.5]{EventID} (itt.north);
\draw[wire] (itt.south) to[out=-70,in=100] node[above left,font=\scriptsize,pos=0.55]{CollectionID} (ct.north);
\draw[wire] (itt.south east) to[out=-50,in=95] node[above right,font=\scriptsize,pos=0.5]{vPEID} (vpet.north);
\draw[wire] (ct.south) -- node[right,font=\scriptsize]{pINTID} (pend.north);
\draw[wire] (vpet.south) -- node[left,font=\scriptsize]{vINTID} (vpend.north);
% direct LPI 旁路：从 Interrupt Source 左缘出发，绕开 Device table 与 Collection table
\draw[dwire] (src.west) to[out=-170,in=95] (1.0,0.0) to[out=-85,in=155] (pend.north west);
\node[left,font=\scriptsize,align=center,text width=1.8cm] at (1.0,0.0) {pINTID (direct LPI support)};
\end{tikzpicture}
\caption[ITS 的 LPI 翻译与路由]{ITS 的 LPI 翻译与路由：DeviceID 经 Device table 选择该设备的 Interrupt translation table，EventID 在 ITT 中翻译为物理路径的 pINTID/CollectionID（经 Collection table 选目标 Redistributor，置 Pending table）或 GICv4 路径的 vINTID/vPEID（经 vPE table 选目标 Redistributor，置 Virtual pending table）；另有不经 ITS 的直接 LPI 旁路。重绘自 Arm GIC Architecture Specification（IHI 0069D）Figure 6-1\cite{source9}。}
\label{fig:gic-its}
\end{figure}
